\section{Variance scaling across practical preprocessing arms}
\label{sec:variance-scaling-followup}

In a separately declared follow-up, we add variance scaling, $s_j = v_j$
($\delta = 1$), to every preprocessing arm of the sampled binary and Mushroom
studies. We reuse all eight binary laws with 128 replications per law, both
Mushroom sample sizes with 32 partitions per size, and the original training,
validation, and test samples. We use population variances only in the synthetic
arm with known preprocessing. All other arms use training variances with
denominator $n$. We keep the same treatment of constant columns and unknown
categories, as well as the same centering, initialization, and intercept rules.

We give no scaling, SD scaling, damping $.01$, and variance scaling the same
25 rate factors relative to empirical curvature. These comprise the earlier
17-factor Mushroom grid, two smaller factors $10^{-5}$ and $10^{-4}$, and six
successive doublings from $12.8$ through $409.6$. With two batch sizes, each
rule has nominal cost $50M$ processed observations per preprocessing arm and
replication or partition at cap $M$. We jointly select the rate, batch size, and stopping
count by aggregate or rare validation loss. The primary search uses the original
checkpoints, while a secondary search adds intermediate checkpoints. We also
retain the nested five-factor practical and 17-factor Mushroom searches to
compare results at the earlier tuning effort. Test outcomes enter neither
selection nor the choice of grid. We retain initialization, fallbacks for empty
groups, numerical failures, and extreme finite selections.

The motivation for variance scaling under the lasso concerns penalty weights on
the original coefficient scale. Here, we use no explicit penalty. The SGD
multiplier for an active coordinate is $\eta/s_j^2 = \eta/v_j^2$, and
calibrating by empirical curvature also changes $\eta$. We therefore compare
how the rules affect relative learning, noise, and stability under a finite
budget. The motivation from shrinkage does not establish which rule should have
lower prediction risk after tuning.

\paragraph{Mushroom results and search boundaries.}

At $n = 2048$ and $M = 8192$, aggregate selection with the original stops
gives variance scaling aggregate loss $.311988$ (MCSE $.201772$) and rare
loss $.051586$ ($.008319$) when predictors are centered and the intercept is
fitted. Its rare-loss differences relative to none, SD, and damping are
$.045630$ ($.008323$), $.046822$ ($.008189$), and $.050904$ ($.008322$),
respectively. Each exceeds two paired MCSE, while the aggregate differences do
not. For variance scaling, aggregate loss has median $.069802$ and maximum
$6.505207$. Selecting by rare validation raises its aggregate mean to $25.5720$
(MCSE $8.0453$).

For the two primary selections at factor $6.4$, the newly tested neighbor
$12.8$ gives finite but worse validation scores. Across both Mushroom sample
sizes, all preprocessing arms, caps, objectives, and both expanded stopping
searches, no selection with a positive budget reaches either grid boundary.
With the original stops, every selected predictor matches the nested 17-factor
search. Sixty candidate indices change because the deterministic tie rule
assigns the same initialization a lower factor. The refined stopping search
selects some added rates, up to $102.4$, and we report it separately.

\begin{table}[ht]
  \centering
  \begin{tabular}{rlrrrr}
    \toprule
     $n$ & Preprocessing      &    None &      SD &  Damped & Variance \\
    \midrule
     256 & Centered, fixed    & .101062 & .075937 & .077910 &  .079580 \\
     256 & Centered, fitted   & .098799 & .076035 & .078161 &  .079579 \\
     256 & Uncentered, fitted & .136151 & .076391 & .082727 &  .079542 \\
    2048 & Centered, fixed    & .005647 & .004722 & .000652 &  .051601 \\
    2048 & Centered, fitted   & .005955 & .004764 & .000682 &  .051586 \\
    2048 & Uncentered, fitted & .043284 & .005060 & .007037 &  .049465 \\
    \bottomrule
  \end{tabular}
  \caption{
    Mean Mushroom rare-group loss under aggregate validation, cap $8192$, and
    the original stops, using the same 25-factor search for every rule. The
    artifacts report both risks, paired uncertainty, and extreme outcomes for
    every preprocessing arm.
  }
\end{table}

At $n = 256$ with centered predictors and a fitted intercept, variance scaling
increases aggregate loss relative to none by $.007047$ (MCSE $.000555$). The
rare-loss difference of $-.019220$ ($.012351$) is unresolved. Five training
partitions and thirteen validation partitions lack green observations, and we
retain all of them. With uncentered predictors, variance scaling gives aggregate
means $.179701$ (MCSE $.150890$) and $13.787223$ ($12.294085$) at the two
sample sizes. These MCSEs describe variation across partitions of one fixed,
previously studied dataset.

\paragraph{Extreme synthetic outcomes and unresolved cases.}

We evaluate all eight laws and five preprocessing arms at equal effort. For the
rarer C2 law with centered predictors and a fitted intercept, aggregate
selection at cap $2048$ gives variance scaling an exact rare-risk median of
$.6836$, mean $1.727 \times 10^{30}$, MCSE $1.552 \times 10^{30}$, and
maximum $1.974 \times 10^{32}$. Extreme outcomes dominate these means,
including those for SD and damping, so they do not precisely estimate expected
risk. Empty rare validation groups can still conceal severe prediction errors.
For the rare-only signal, variance scaling gives aggregate mean $.01065$
($.00311$) but rare mean $.4122$ ($.1467$), compared with $.2326$ ($.0319$)
for no scaling.

Across the two expanded stopping searches, synthetic selections reach the lower
rate boundary 106 times and the upper boundary 52 times, counting repeated
queries on the same replications. Only variance scaling reaches the upper
boundary, under known preprocessing in the two rarer anchor laws. Resolving the
Mushroom boundaries therefore leaves these synthetic cases open. Across both
stages, the archives retain 6,451,200 endpoints and 387,072 choices, including
1,031,540 failed endpoints on 330,918 failed paths.

\paragraph{Independent audit and numerical limits.}

The audit reconstructs every loss and selection, verifies all 179 preserved
files, and replays 100,271 paths in explicitly transformed coordinates.
Overlapping frozen coefficient paths reproduce exactly. Reconstructed losses
have maximum normwise discrepancy $1.36 \times 10^{-14}$ and leave every
selection unchanged. Mushroom passes the declared trajectory and prediction
tolerances. For synthetic data, however, the audit retains 63 discrepancies in
twelve replications, all of which have empty rare validation groups. These
comprise 60 checks of selected predictions, two first-failure times, and one
trajectory check. Repeated queries can refer to the same selected path. The
maximum normwise coefficient and prediction discrepancies are $.003747$ and
$.004184$, exceeding the declared tolerances of $2 \times 10^{-9}$ and
$10^{-8}$. All stored checkpoint statuses agree.

The transformed replay therefore does not certify these extreme synthetic
paths. We retain the original candidates and selections in the archive. The
result note records the affected configurations and explains how we extended the
checker to retain disagreements in failure times rather than stop at the first
one. See \path{notes/variance-scaling-followup-results.md} for the protocol,
archives, and reproduction commands. All earlier studies retain their frozen
file hashes.
